% Bagean: metodhe
% Bab: bukti
% Pérangan: conto-1

\documentclass[../../../include/open-logic-section]{subfiles}

\begin{document}

\olfileid{mth}{prf}{ex1}

\olsection{Sawijining Conto}

Conto kapisan yaiku kasunyatan prasaja ing ngisor iki babagan
gabungan lan irisan himpunan. Conto iki bakal nuduhaké cara
njlentrehaké dhéfinisi, mbuktekaké konjungsi lan pratelan
universal, sarta mbuktekaké kanthi kasus.

\begin{prop}
Kanggo sembarang himpunan $A$, $B$, lan $C$, $A \cup (B \cap C) = (A \cup B)
\cap (A \cup C)$
\end{prop}

Ayo dibuktèkaké!{}

\begin{proof}
Kita arep nuduhaké yèn kanggo sembarang himpunan $A$, $B$, lan $C$,
$A \cup (B \cap C) = (A \cup B) \cap (A \cup C)$.
\begin{quote}
Kaping pisan, kita njlentrehaké dhéfinisi ``$=$'' ing pratelan
proposisi. Élinga yèn mbuktekaké pepadhan himpunan ateges
mbuktekaké yèn loro himpunan nduwèni !!{element}s kang padha.
Tegesé, kabèh !!{element}s saka $A \cup (B \cap C)$ uga
!!{element}s saka $(A \cup B) \cap (A \cup C)$, lan kosok baliné.
``Kosok baliné'' ateges saben !!{element} saka $(A \cup B)
\cap (A \cup C)$ uga kudu !!a{element} saka $A \cup (B \cap
C)$. Dadi, sawisé dhéfinisi dijlentrehaké, katon yèn kita
kudu mbuktekaké konjungsi. Ayo iki dicathet:
\end{quote}
Miturut dhéfinisi, $A \cup (B \cap C) = (A \cup B) \cap (A \cup C)$
yèn lan mung yèn saben !!{element} saka $A \cup (B \cap C)$
uga !!a{element} saka $(A \cup B) \cap (A \cup C)$, lan
saben !!{element} saka $(A \cup B) \cap (A \cup C)$ uga
!!a{element} saka $A \cup (B \cap C)$.
\begin{quote}
Amarga iki konjungsi, saben pérangan kudu dibuktèkaké dhéwé.
Kita wiwiti saka kang kapisan: buktèkna yèn saben
!!{element} saka $A \cup (B \cap C)$ uga !!a{element}
saka $(A \cup B) \cap (A \cup C)$.

Iki pratelan universal. Mula kita jupuk !!{element} sembarang
saka $A \cup (B \cap C)$, banjur nuduhaké yèn iku uga
!!a{element} saka $(A \cup B) \cap (A \cup C)$.
Kita wènèhi jeneng variabel kanggo !!{element} sembarang
iku, upamané~$z$. Buktiné diterusaké:
\end{quote}
Kaping pisan, kita buktèkaké yèn saben !!{element} saka
$A \cup (B \cap C)$ uga !!a{element} saka $(A \cup B)
\cap (A \cup C)$. Jupuken $z \in A \cup (B
\cap C)$. Kita kudu nuduhaké yèn
$z \in (A \cup B) \cap (A \cup C)$.
\begin{quote}
Saiki dhéfinisi $\cup$ lan~$\cap$ kudu dijlentrehaké.
Upamané, dhéfinisi $\cup$ yaiku
$A \cup B = \Setabs{z}{z \in A
  \text{ utawa } z \in B}$. Nalika dhéfinisi iki ditrapaké
marang ``$A \cup (B
\cap C)$'', panggonan ``$B$'' ing dhéfinisi
saiki diisi déning ``$B \cap C$'', saéngga
$A \cup (B \cap C) = \Setabs{z}{z \in A \text{ utawa }
  z \in B \cap C}$. Dadi anggepan
$z \in A \cup (B \cap C)$ padha karo
$z \in \Setabs{z}{z \in A \text{ utawa } z \in B \cap
  C}$. Lan $z \in \Setabs{z}{\dots z\dots}$ yèn lan mung yèn
\dots $z$ \dots; ing kéné tegesé $z \in A$ utawa $z \in B \cap C$.
\end{quote}
Miturut dhéfinisi $\cup$, salah siji $z \in A$ utawa
$z \in B \cap C$.
\begin{quote}
Amarga iki disjungsi, bukti mawa kasus migunani. Kita
mbédakaké rong kasus lan nuduhaké yèn saben kasus ngasilaké
kesimpulan kang dituju, yaiku ``$z \in (A \cup B) \cap (A \cup
C)$''.
\end{quote}
Kasus 1: Anggepen $z \in A$.
\begin{quote}
Anggepan iki ora akèh mènèhi bahan tambahan. Mula kita deleng
dhisik wujud kesimpulan: kita arep nuduhaké
$z \in (A \cup B) \cap (A \cup C)$. Miturut dhéfinisi
$\cap$, yèn arep nuduhaké $z \in (A \cup B) \cap (A \cup C)$,
kita kudu nuduhaké yèn obyèk mau dadi anggota $(A \cup B)$
lan $(A \cup C)$. Nanging $z
\in A \cup B$ yèn lan mung yèn $z \in A$ utawa $z \in B$,
lan anggepan kasus~1 wis mènèhi $z \in A$.
Kanthi panalaran padha, kanthi ngganti $C$ kanggo $B$,
kita olèh $z \in A \cup C$.
Panalaran mau mlaku saka arah kosok, mula saiki ayo
kita tulis manut urutan kang dibutuhaké bukti.
\end{quote}
Amarga $z \in A$, kita olèh $z \in A$ utawa $z \in B$;
mula miturut dhéfinisi~$\cup$, $z \in A \cup B$.
Semono uga $z \in A \cup C$. Iki ateges
$z \in (A \cup B) \cap (A \cup C)$ miturut dhéfinisi~$\cap$.
\begin{quote}
Iki ngrampungaké kasus kapisan. Saiki kesimpulan padha
kudu diturunaké ing kasus kapindho, yaiku nalika
$z \in B \cap C$.
\end{quote}
Kasus 2: Anggepen $z \in B \cap C$.
\begin{quote}
Kita ketemu maneh karo irisan rong himpunan.
Ayo dhéfinisi~$\cap$ ditrapaké:
\end{quote}
Amarga $z \in B \cap C$, $z$ mesthi !!a{element} saka
$B$ lan $C$, miturut dhéfinisi~$\cap$.
\begin{quote}
Saiki delengen kesimpulan maneh. Kita kudu nuduhaké yèn
$z$ ana ing $(A \cup B)$ lan $(A \cup C)$.
Wangsulané uga langsung.
\end{quote}
Amarga $z \in B$, kita olèh $z \in (A \cup B)$.
Amarga $z \in C$, uga $z \in (A
\cup C)$. Dadi $z \in (A \cup B) \cap (A \cup C)$.
\begin{quote}
Ing kéné dhéfinisi $\cup$ lan $\cap$ ditrapaké maneh.
Amarga dhéfinisi kasebut wis dielingaké lan kita wis
nuduhaké yèn $z$ kang ana ing salah siji saka rong
himpunan uga ana ing gabungané, langkah iki ora perlu
dijlentrehaké luwih rinci.

Kasus kapindho uga rampung; mula kesimpulan kapisan
bisa dipratelakaké.
\end{quote}
Dadi, yèn $z \in A \cup (B \cap C)$ mula
$z \in (A \cup B) \cap (A \cup C)$.
\begin{quote}
Saiki kari mbuktekaké arah kosok baliné: saben
!!{element} saka $(A \cup B) \cap (A \cup C)$
uga !!a{element} saka $A \cup (B \cap
C)$. Kaya sadurungé, pratelan universal iki dibuktèkaké
kanthi njupuk !!{element} sembarang saka himpunan
kapisan lan nuduhaké yèn iku anggota himpunan kapindho.
Ayo pratelakna tujuwan iki.
\end{quote}
Saiki anggepen $z \in (A \cup B) \cap (A \cup C)$.
Kita arep nuduhaké yèn $z \in A \cup (B \cap C)$.
\begin{quote}
Saiki kita mangkat saka anggepan $z \in (A \cup B) \cap (A
\cup C)$. Panganggoné~$z$ kang padha kaya ing pérangan
kapisan muga-muga ora mbingungaké. Sawisé pérangan iku
rampung, kabèh anggepan ing kono wis ora lumaku, mula
saiki kita bisa nggawe anggepan anyar babagan $z$.
Yèn iki mbingungaké, gantinen $z$ nganggo variabel liya
ing bukti sabanjuré.

Kita ngerti yèn $z$ ana ing $A \cup B$ lan
$A \cup C$, miturut dhéfinisi~$\cap$. Miturut dhéfinisi
$\cup$, iki bisa dijlentrehaké maneh: $z \in A$
utawa $z \in B$, lan uga $z \in A$ utawa $z \in
C$. Iki kaya bukti mawa kasus, nanging tembung ``lan''
bisa mbingungaké. Mbokmenawa panjenengan ngira yèn
iki mung ana telung kamungkinan: $z$ ana ing $A$, $B$,
utawa $C$. Iku klèru. Kita kudu ngati-ati lan nliti
saben disjungsi siji-siji.
\end{quote}
Miturut dhéfinisi $\cap$, $z \in A \cup B$ lan
$z \in A \cup C$. Miturut dhéfinisi $\cup$,
$z \in A$ utawa $z \in B$. Kita mbédakaké kasus.
\begin{quote}
Amarga fokus ing disjungsi kapisan, disjungsi kapindho
saka njlentrehaké $A \cup C$ durung ditulis.
Saiki pancèn durung perlu. Kasus kapisan yaiku
$z \in A$, lan !!a{element} sawijining himpunan
uga !!a{element} gabungané karo himpunan liya.
Dadi kasus~1 gampang:
\end{quote}
Kasus 1: Anggepen $z \in A$. Mula
$z \in A \cup (B \cap
C)$.
\begin{quote}
Saiki kasus kapindho yaiku $z \in B$.
Ing kéné kita njlentrehaké $\cup$ kapindho lan
mbédakaké kasus maneh:
\end{quote}
Kasus 2: Anggepen $z \in B$. Amarga $z \in A \cup C$,
$z \in A$ utawa $z \in C$. Kita mbédakaké kasus maneh:

Kasus 2a: $z \in A$. Mula maneh,
$z \in A \cup (B \cap C)$.
\begin{quote}
Iki rada anèh. Anggepan yèn~$z \in B$ pancèn ora
dibutuhaké ing kasus iki, nanging ora dadi masalah.
\end{quote}
Kasus 2b: $z \in C$. Mula $z \in B$ lan $z \in C$,
saéngga $z \in B \cap C$; banjur
$z \in A \cup (B \cap C)$.
\begin{quote}
Iki ngrampungaké loro bukti mawa kasus, lan kanthi
iku pérangan kapindho uga rampung.
\end{quote}
Dadi, yèn $z \in (A \cup B) \cap (A \cup C)$ mula
$z \in A \cup (B \cap C)$.
\end{proof}

\end{document}
